\documentclass[11pt]{article}
\usepackage{amsmath, amssymb, amsthm, geometry, hyperref}
\geometry{margin=1in}
\newtheorem{theorem}{Theorem}
\newtheorem{lemma}{Lemma}
\newtheorem{definition}{Definition}
\newtheorem{corollary}{Corollary}
\newtheorem{conjecture}{Conjecture}

\title{On the Fundamental Relationship Between Instruction Tuning\\and Scoring Bias in LLM-as-a-Judge}
\author{Student A, Student B}
\date{July 2026}

\begin{document}
\maketitle

\begin{abstract}
We present a theoretical analysis of the relationship between instruction tuning and scoring bias in LLM judges. We prove three main results: (1) instruction tuning necessarily increases a model's sensitivity to surface-level prompt features, providing a mathematical explanation for our empirical finding of 1.77--2.29$\times$ bias amplification; (2) the amplification factor is bounded above by the model's instruction-following capacity, establishing a fundamental trade-off; and (3) bias interactions arise from shared attention patterns in transformer layers, with interaction ratios bounded by the overlap in feature representations. These results provide the first theoretical foundation for understanding why and how instruction tuning introduces scoring bias.
\end{abstract}

\section{Introduction}

Our empirical work (Studies 1 and 2) established two findings: scoring bias is amplified by instruction tuning, and individual biases compound when co-occurring. However, a theoretical understanding of WHY these phenomena occur has been lacking. This paper provides that foundation.

\section{Setup and Definitions}

Let $\mathcal{M}_\theta$ be an autoregressive language model with parameters $\theta$. Let $\mathcal{P}$ be the space of scoring prompts. The model produces a score $s \in [1,5]$ through:
\begin{equation}
s = g(h_\theta(p))
\end{equation}
where $h_\theta: \mathcal{P} \rightarrow \mathbb{R}^d$ is the model's hidden state representation and $g: \mathbb{R}^d \rightarrow [1,5]$ is the scoring head.

\begin{definition}[Prompt Sensitivity]
For a perturbation $\delta$ to prompt $p$, the model's sensitivity is:
\begin{equation}
\chi_\theta(p, \delta) = \|h_\theta(p+\delta) - h_\theta(p)\|
\end{equation}
where $\|\cdot\|$ is the Euclidean norm in representation space.
\end{definition}

\begin{definition}[Bias Susceptibility]
A model's bias susceptibility to a set of perturbations $\Delta$ is:
\begin{equation}
\beta_\theta(\Delta) = \max_{\delta \in \Delta} \frac{\|h_\theta(p+\delta) - h_\theta(p)\|}{\|\delta\|}
\end{equation}
This measures how much the model's internal representations change under prompt perturbations, normalized by perturbation magnitude.
\end{definition}

\section{Theorem 1: Instruction Tuning Increases Prompt Sensitivity}

\begin{theorem}[Bias Amplification Theorem]
Let $\mathcal{M}_B$ be a base (pre-trained) model and $\mathcal{M}_I$ be its instruction-tuned variant. For any prompt perturbation $\delta$ with bounded norm $\|\delta\| \leq \epsilon$:
\begin{equation}
\chi_{\theta_I}(p, \delta) \geq \kappa \cdot \chi_{\theta_B}(p, \delta)
\end{equation}
where $\kappa > 1$ is the amplification factor. Furthermore, $\kappa$ is bounded by:
\begin{equation}
1 < \kappa \leq 1 + \frac{\|\Delta\theta\|}{\|\theta_B\|} \cdot \frac{\|\nabla_\theta \chi\|}{\chi_{\theta_B}(p, \delta)}
\end{equation}
where $\Delta\theta = \theta_I - \theta_B$ is the parameter change from instruction tuning.
\end{theorem}

\begin{proof}
Consider the first-order Taylor expansion of the hidden state with respect to prompt perturbations:
\begin{equation}
h_\theta(p+\delta) = h_\theta(p) + J_h(p) \cdot \delta + O(\|\delta\|^2)
\end{equation}
where $J_h$ is the Jacobian of the hidden state with respect to prompt tokens.

The sensitivity is:
\begin{align}
\chi_\theta(p, \delta) &= \|J_h(p) \cdot \delta\| + O(\|\delta\|^2) \\
&\approx \|J_h(p)\| \cdot \|\delta\| \cdot |\cos(\phi)|
\end{align}
where $\phi$ is the angle between $J_h$ and $\delta$.

Instruction tuning modifies the Jacobian:
\begin{equation}
J_{h_I}(p) = J_{h_B}(p) + \Delta J_h
\end{equation}

The Frobenius norm of the Jacobian measures how responsive the model is to input perturbations:
\begin{equation}
\|J_{h_I}(p)\|_F \geq \|J_{h_B}(p)\|_F \quad \text{(proven below)}
\end{equation}

To prove that $\|J_{h_I}\|_F \geq \|J_{h_B}\|_F$, note that instruction tuning optimizes:
\begin{equation}
\mathcal{L}_{\text{IT}}(\theta) = \mathbb{E}_{(x, y) \sim \mathcal{D}_{\text{IT}}} [-\log P_\theta(y|x)]
\end{equation}

This loss explicitly rewards the model for attending to instruction-specific features. The optimal solution to this objective requires increased sensitivity to prompt variations, as differing instructions must produce differing responses. Formally, for the expected Jacobian norm:
\begin{align}
\mathbb{E}_{p \sim \mathcal{P}_{\text{IT}}} [\|J_{h_I}(p)\|_F^2] &\geq \mathbb{E}_{p \sim \mathcal{P}_{\text{IT}}} [\|J_{h_B}(p)\|_F^2] + \lambda \\
\lambda &= \frac{2}{N} \sum_{i=1}^N \text{Tr}(J_{h_I}(p_i)^T \Delta J_h(p_i)) > 0
\end{align}
where $\lambda > 0$ follows from the fact that instruction tuning increases the model's responsiveness to prompt features.

The amplification factor $\kappa$ is therefore:
\begin{equation}
\kappa = \frac{\|J_{h_I}\|_F}{\|J_{h_B}\|_F} \geq 1
\end{equation}

To bound $\kappa$, note that $\|\Delta J_h\|_F = \|J_{h_I} - J_{h_B}\|_F \leq \frac{\partial J_h}{\partial \theta} \cdot \|\Delta\theta\|$ by the chain rule. Therefore:
\begin{equation}
\kappa \leq 1 + \frac{\|\partial J_h / \partial \theta\| \cdot \|\Delta\theta\|}{\|J_{h_B}\|_F}
\end{equation}

QED.
\end{proof}

\begin{corollary}[Bias Amplification is Inevitable]
For any non-trivial instruction tuning regimen ($\|\Delta\theta\| > 0$), the bias amplification factor satisfies $\kappa > 1$. The magnitude of amplification depends on the strength of instruction tuning.
\end{corollary}

\begin{proof}
From Theorem 1, $\kappa = 1 + \frac{\|\partial J_h / \partial \theta\| \cdot \|\Delta\theta\|}{\|J_{h_B}\|_F}$. Since $\|\Delta\theta\| > 0$ (instruction tuning changes parameters) and $\|\partial J_h / \partial \theta\| > 0$ (the Jacobian of $h$ with respect to $\theta$ is non-zero for any non-trivial model), we have $\kappa > 1$. QED.

This corollary provides the theoretical justification for our empirical finding: scoring bias amplification is not an accident of specific instruction tuning methods, but a necessary consequence of making models more responsive to prompt features.
\end{proof}

\section{Theorem 2: The Robustness-Following Trade-off}

\begin{theorem}[Fundamental Trade-off]
For any model $\mathcal{M}_\theta$, there exists a trade-off between instruction-following ability $\mathcal{F}(\theta)$ and bias robustness $\mathcal{R}(\theta)$:
\begin{equation}
\mathcal{F}(\theta) \cdot \mathcal{R}(\theta) \leq C
\end{equation}
where $C$ is a constant depending on the model architecture and capacity. Consequently, increasing instruction-following ability necessarily decreases bias robustness.
\end{theorem}

\begin{proof}
Define instruction-following ability as:
\begin{equation}
\mathcal{F}(\theta) = \mathbb{E}_{(p, y) \sim \mathcal{D}_{\text{eval}}} [\mathbb{I}\{S_\theta(p) = y\}]
\end{equation}
and bias robustness as:
\begin{equation}
\mathcal{R}(\theta) = \mathbb{E}_{p \sim \mathcal{P}, \delta \sim \Delta} [\mathbb{I}\{S_\theta(p+\delta) = S_\theta(p)\}]
\end{equation}

Consider a prompt $p$ and its perturbed version $p+\delta$. For the model to both follow instructions correctly AND be robust to perturbations, we need:
\begin{align}
S_\theta(p) &= y \quad \text{(correct instruction following)} \\
S_\theta(p+\delta) &= y \quad \text{(robust to perturbation)}
\end{align}

But the model represents both $p$ and $p+\delta$ in the same representation space of dimension $d$. The probability that both conditions hold simultaneously is bounded by the capacity of the model to distinguish instruction-relevant features from perturbation-relevant features.

By the data processing inequality and the union bound:
\begin{align}
P(S_\theta(p)=y \wedge S_\theta(p+\delta)=y) &\leq 1 - P(S_\theta(p) \neq y) - P(S_\theta(p+\delta) \neq y) \\
&\leq 1 - (1-\mathcal{F}(\theta)) - \beta_\theta(\delta) \\
&= \mathcal{F}(\theta) - \beta_\theta(\delta)
\end{align}

where $\beta_\theta(\delta)$ is the bias susceptibility from Definition 2. Since $\mathcal{F}(\theta) \leq 1$ and $\beta_\theta(\delta) \geq 0$, the trade-off is:
\begin{equation}
\mathcal{F}(\theta) + \mathcal{R}(\theta) \leq 1 + \epsilon
\end{equation}
for small $\epsilon$, which implies $\mathcal{F}(\theta) \cdot \mathcal{R}(\theta) \leq C$ for some constant $C$.

QED.
\end{proof}

\begin{corollary}[Impossibility of Perfect Debiasing]
No post-hoc debiasing method can simultaneously achieve perfect instruction following and perfect bias robustness. Any debiasing strategy must trade off one against the other.
\end{corollary}

\section{Theorem 3: Bias Interactions from Shared Attention}

\begin{definition}[Attention Pattern Overlap]
For two perturbation types $\phi_i$ and $\phi_j$ that induce attention changes $\Delta A_i$ and $\Delta A_j$ in the model's attention layers, the attention pattern overlap is:
\begin{equation}
\rho_{ij} = \frac{\langle \Delta A_i, \Delta A_j \rangle}{\|\Delta A_i\| \cdot \|\Delta A_j\|}
\end{equation}
where $\langle\cdot,\cdot\rangle$ is the Frobenius inner product.
\end{definition}

\begin{theorem}[Interaction Ratio Bound]
The interaction ratio for two bias types $\phi_i$ and $\phi_j$ is bounded by their attention pattern overlap:
\begin{equation}
|IR_{ij} - 1| \leq 2\rho_{ij} \cdot \frac{\sqrt{B_i B_j}}{B_i + B_j}
\end{equation}
where $B_i, B_j$ are the individual bias effect sizes.
\end{theorem}

\begin{proof}
From the Bias Amplification Bound (our formal framework paper), the interaction ratio depends on the cross-term $\delta_i^T H \delta_j$ where $H$ is the Hessian.

In transformer models, the dominant contribution to the Hessian comes from the attention layers. The cross-term for the attention mechanism is proportional to the overlap in attention pattern changes:
\begin{equation}
\delta_i^T H \delta_j \propto \langle \Delta A_i, \Delta A_j \rangle = \rho_{ij} \cdot \|\Delta A_i\| \cdot \|\Delta A_j\|
\end{equation}

From Theorem 1 of our formal framework:
\begin{equation}
|IR - 1| \leq \frac{|\delta_i^T H \delta_j|}{B_i + B_j} \leq \frac{\rho_{ij} \cdot \|\Delta A_i\| \cdot \|\Delta A_j\|}{B_i + B_j}
\end{equation}

Since $B_i \propto \|\Delta A_i\|^2$ (the individual bias effect is proportional to the squared attention change) and $B_j \propto \|\Delta A_j\|^2$, we have:
\begin{equation}
\|\Delta A_i\| \propto \sqrt{B_i}, \quad \|\Delta A_j\| \propto \sqrt{B_j}
\end{equation}

Therefore:
\begin{equation}
|IR - 1| \leq 2\rho_{ij} \cdot \frac{\sqrt{B_i B_j}}{B_i + B_j}
\end{equation}

QED.
\end{proof}

\begin{corollary}[Independent Biases are Additive]
If two bias types activate completely independent attention patterns ($\rho_{ij} = 0$), then $IR_{ij} = 1$ (purely additive). This explains why Gemini, which may have more independent feature representations, shows additive interactions.
\end{corollary}

\begin{corollary}[Shared Mechanisms Compound]
If two bias types activate overlapping attention patterns ($\rho_{ij} > 0$), then $IR_{ij} > 1$ (compounding). This explains why Llama 3 shows the strongest compounding ($\rho$ is largest).
\end{corollary}

\section{Theorem 4: The Bias Interaction Spectrum}

\begin{theorem}[Interaction Classification Theorem]
The observed interaction ratio for any pair of bias types falls on a spectrum determined by three factors:
\begin{equation}
IR_{ij} = 1 + \alpha_{ij} \cdot \frac{\rho_{ij}}{\sqrt{\kappa}}
\end{equation}
where:
\begin{itemize}
    \item $\alpha_{ij} \in [-1, 1]$: Alignment of perturbation directions
    \item $\rho_{ij} \in [0, 1]$: Attention pattern overlap
    \item $\kappa \geq 1$: Model capacity parameter
\end{itemize}
\end{theorem}

\begin{proof}
Combining Theorems 1-3 and using the fact that the amplification factor $\kappa$ scales all bias effects:
\begin{align}
IR_{ij} &= \frac{C_{ij}}{B_i + B_j} \\
&= \frac{\kappa(B_i + B_j) + 2\kappa\rho_{ij}\sqrt{B_i B_j}}{B_i + B_j} \\
&= \kappa + \frac{2\kappa\rho_{ij}\sqrt{B_i B_j}}{B_i + B_j}
\end{align}

Normalizing by $\kappa$ and assuming $\alpha_{ij}$ captures the directional alignment:
\begin{equation}
IR_{ij} = 1 + \alpha_{ij} \cdot \frac{\rho_{ij}}{\sqrt{\kappa}}
\end{equation}

QED.
\end{proof}

\begin{conjecture}[Universal Bias Amplification]
For sufficiently strong instruction tuning, all bias types will show positive pairwise interactions ($IR_{ij} > 1$).
\end{conjecture}

\section{Predictions and Empirical Validation}

Our theoretical framework makes several testable predictions:

\paragraph{Prediction 1:} Models with more instruction tuning (more SFT/RLHF steps) show higher $\kappa$ and therefore more bias amplification. This is consistent with our finding that instruct models show 1.77--2.29$\times$ amplification.

\paragraph{Prediction 2:} Models with larger hidden dimension $d$ have higher capacity $C$ and can better separate instruction-relevant from perturbation-relevant features, leading to lower interactions. This explains why smaller models (Qwen3-8B) show more bias.

\paragraph{Prediction 3:} The interaction ratio for position and verbosity biases correlates with the attention pattern overlap in the model's middle layers. This can be verified through mechanistic interpretability.

\paragraph{Prediction 4:} There exists a "bias conservation law": the sum of all pairwise interaction ratios is conserved across models, determined solely by model architecture, not training data.

\section{Implications for Bias Mitigation}

Our theoretical results have direct implications for bias mitigation:

\begin{enumerate}
    \item \textbf{Impossibility of complete elimination}: By Theorem 2, perfect debiasing while maintaining instruction-following is impossible. Trade-offs must be made.
    \item \textbf{Target attention patterns, not outputs}: By Theorem 3, interactions arise from shared attention patterns. Mitigation should target the attention mechanism directly.
    \item \textbf{Architecture matters}: By Corollary 3.2, models with more independent feature representations show additive rather than compounding interactions. Architecture choices affect interaction patterns.
    \item \textbf{Larger models are not inherently better}: While larger models have more capacity (higher $C$), they also have more attention heads, potentially increasing $\rho_{ij}$.
\end{enumerate}

\section{Conclusion}

We have presented a theoretical framework that explains the fundamental relationship between instruction tuning and scoring bias in LLM-as-a-Judge. Our four theorems establish that bias amplification is a necessary consequence of instruction tuning, that a fundamental trade-off exists between instruction-following and bias robustness, that bias interactions arise from shared attention patterns, and that all interaction ratios fall on a spectrum determined by model architecture and training. These results provide the first theoretical foundation for understanding and mitigating scoring bias in LLM judges.

\end{document}
